% Put your abstract here

\begin{abstractwithpageno}	
Agriculture remains a primary source of livelihood in the Municipality of Sinacaban, Misamis Occidental, where many farmers rely on seasonal crops. Interviews with local farmers and the Municipal Agriculture Office revealed recurring problems such as lack of timely price information, weak linkages with buyers, delays in transactions, and limited visibility of available produce at the barangay level. These challenges contribute to low profitability and underutilized harvests.

This project develops the Digital Farmers’ Produce System, a cloud-based web application designed to help farmers manage their seasonal crop listings, request support from the local agriculture office, and connect with potential buyers. The system allows farmers to publish their available produce with quantity and harvest schedules, buyers to search and view listings by location and category, and agriculture officers to monitor barangay-level stock information and resource requests.

The system was developed using the Rapid Application Development (RAD) framework and implemented using PHP 8, the Flight microframework, MySQL, HTML, CSS, JavaScript, Bootstrap, and shared cloud hosting. Security and reliability considerations included role-based access, encrypted communication through HTTPS, input validation, prepared statements, password hashing, and backup planning. User acceptance was evaluated with 12 respondents: five farmers, four buyers, and three DA/MAO officers. Results showed strongly favorable ratings across role groups, with overall means of 4.60 for farmers, 4.83 for buyers, and 4.69 for DA/MAO officers on a 5-point scale. These results suggest that a localized, user-centered agricultural platform can strengthen market visibility and institutional coordination, although wider deployment will require continued training, connectivity support, and security validation before production rollout.
\end{abstractwithpageno}